\documentclass[11pt,a4paper]{article}

\usepackage[margin=1.05in]{geometry}
\usepackage{amsmath,amssymb,amsthm}
\usepackage{mathtools}
\usepackage{stmaryrd}
\usepackage{microtype}
\setlength{\emergencystretch}{2.5em}
\usepackage{enumitem}
\usepackage{booktabs}
\usepackage{array}
\usepackage{tikz-cd}
\usepackage{xcolor}
\usepackage[colorlinks=true,linkcolor=blue!45!black,urlcolor=blue!60!black,citecolor=blue!45!black]{hyperref}

\newtheorem{theorem}{Theorem}[section]
\newtheorem{proposition}[theorem]{Proposition}
\newtheorem{lemma}[theorem]{Lemma}
\newtheorem{corollary}[theorem]{Corollary}
\newtheorem{conjecture}[theorem]{Conjecture}
\theoremstyle{definition}
\newtheorem{definition}[theorem]{Definition}
\newtheorem{example}[theorem]{Example}
\theoremstyle{remark}
\newtheorem{remark}[theorem]{Remark}
\newtheorem{warning}[theorem]{Warning}
\newtheorem{question}[theorem]{Question}

\newcommand{\FF}{\mathbb{F}}
\newcommand{\NN}{\mathbb{N}}
\newcommand{\ZZ}{\mathbb{Z}}
\newcommand{\QQ}{\mathbb{Q}}
\newcommand{\cat}[1]{\mathbf{#1}}
\newcommand{\Set}{\cat{Set}}
\newcommand{\Grpd}{\cat{Grpd}}
\newcommand{\FinSet}{\cat{FinSet}}
\newcommand{\Vect}{\cat{Vect}}
\newcommand{\vect}{\cat{vect}}
\newcommand{\op}{\mathrm{op}}
\DeclareMathOperator{\Hom}{Hom}
\DeclareMathOperator{\GL}{GL}
\DeclareMathOperator{\PGL}{PGL}
\DeclareMathOperator{\Aut}{Aut}
\DeclareMathOperator{\Ind}{Ind}
\DeclareMathOperator{\Inc}{Inc}
\DeclareMathOperator{\deng}{deg}
\newcommand{\qbin}[2]{\genfrac{[}{]}{0pt}{}{#1}{#2}_{q}}
\newcommand{\norm}[1]{\lvert #1 \rvert}
\newcommand{\iso}{\mathrm{iso}}

\definecolor{okgreen}{rgb}{0.0,0.45,0.15}
\definecolor{nored}{rgb}{0.65,0.0,0.0}
\newcommand{\OK}{\textcolor{okgreen}{\textbf{verified}}}
\newcommand{\NO}{\textcolor{nored}{\textbf{fails}}}

\title{\bfseries Dirichlet Convolution over a Finite Field\\[0.4em]
\large The multiplicative leg of Day convolution, and the torus that obstructs it}
\author{}
\date{}

\begin{document}
\maketitle

\begin{abstract}
Day convolution along $\sqcup$ on finite sets decategorifies to the Cauchy product of
exponential generating functions; Day convolution along $\times$ decategorifies to the
Dirichlet product. Over $\FF_q$ the additive leg is well understood --- it splits into a
Levi version and a parabolic version, and both appear in the literature. The
multiplicative leg does not appear at all. We define it, as Day convolution along
$\otimes$ on the groupoid $\vect_q^{\iso}$ of finite-dimensional $\FF_q$-vector spaces,
and identify its incidence algebra with the \emph{$q$-Dirichlet series}
$\sum_{n\ge1} a_n\, n^{-s}/\norm{\GL_n(\FF_q)}$.

The main structural finding is negative and, we think, the real content of the question.
The additive leg works over $\FF_q$ because the torus factor $(q-1)^{\dim}$ in
$\norm{\GL_n(\FF_q)}$ cancels in a convolution indexed by a \emph{sum} of dimensions. It
does not cancel in a convolution indexed by a \emph{product}, because $ab \neq a+b$. In
consequence the multiplicative leg has no $q \to 1$ limit, its convolution unit is
$(q-1)\delta^1$ rather than $\delta^1$, and the $\FF_1$ dictionary breaks at exactly this
point. Removing the torus by hand repairs the limit and produces a polynomial
$q$-analogue
\[
  D_{a,b}(q) \;=\; q^{\binom{ab}{2}-\binom{a}{2}-\binom{b}{2}}\,
  \frac{[ab]_q!}{[a]_q!\,[b]_q!}
  \;\in\; \NN[q],
  \qquad D_{a,b}(1) = \frac{(ab)!}{a!\,b!},
\]
but this normalisation is not a groupoid cardinality, so the repair costs the objective
interpretation. We preserve the disagreement rather than resolving it.

We contrast this with the other $\FF_q$-analogue of Dirichlet convolution --- the monoid
of monic polynomials in $\FF_q[t]$, or of effective divisors on a curve --- which has an
honest Euler product and an honest zeta function, but places $q$ in the cardinality of
the index monoid rather than in automorphisms. The two are genuinely different
deformations of the two things that $(\NN_{>0},\times)$ does at once, and only the Hall
algebra of $\mathrm{Coh}(X)$ appears to do both.
\end{abstract}

\tableofcontents

\section{The question}

The slogan is that Day convolution is a tensor product is a Dirichlet series. Unpacked:

\begin{definition}[Day convolution \cite{Day}]\label{def:day}
For a small monoidal category $(\mathcal{A},\otimes)$, the functor category
$[\mathcal{A},\Set]$ carries
\[
  (F \star G)(c) \;=\; \int^{a,b} \mathcal{A}(a \otimes b,\, c) \times F(a) \times G(b),
\]
the unique monoidal structure making the Yoneda embedding strong monoidal.
\end{definition}

The point is that $\mathcal{A}$ is an input, not a constant. Vary it and the
decategorified product varies:

\begin{center}
\small
\begin{tabular}{@{}lll@{}}
\toprule
$\mathcal{A}$ & Day convolution & decategorification \\
\midrule
$(\FinSet^{\iso}, \sqcup)$ & species product & $\sum a_n x^n/n!$ (Cauchy/EGF) \\
$(\NN_{>0}, \times)$ & Dirichlet convolution & $\sum a_n n^{-s}$ \\
$(\FinSet^{\iso}_{\neq\emptyset}, \times)$ & arithmetic product & $\sum a_n n^{-s}/n!$ \\
\bottomrule
\end{tabular}
\end{center}

\noindent
The third row is Maia--M\'endez's arithmetic product \cite{MaiaMendez} and
Baez--Dolan's arithmetic species \cite{BaezDolan}; it is the ``fat'' version of the
second. Observe what the first and third rows have in common: $(\FinSet^{\iso},\sqcup,
\times)$ is a \emph{rig} (bimonoidal) category, and the two Day convolutions are the two
monoidal structures on one object.

The question of this note is what happens when $\FinSet$ is replaced by $\vect_q$, the
category of finite-dimensional vector spaces over $\FF_q$. That category is a rig
category too, with $\oplus$ and $\otimes$. So the question has a precise form:

\begin{quote}
\emph{What are the two Day convolutions on $\vect_q$, and which classical series do they
decategorify to?}
\end{quote}

The $\oplus$ leg turns out to be thoroughly studied, under several names, and to split
into two inequivalent versions (\S\ref{sec:known}). The $\otimes$ leg appears to be
absent from the literature. \S\ref{sec:tensor}--\S\ref{sec:torus} define it and explain
why it has been absent: it is obstructed, and the obstruction is the maximal torus.

\section{What is already known}\label{sec:known}

\subsection{The additive leg, twice}

Gálvez-Carrillo, Kock and Tonks \cite{GKT-comb} survey incidence algebras arising from
decomposition spaces (equivalently $2$-Segal spaces \cite{DyckerhoffKapranov}), and their
list of power series representations contains six entries. Two of them are over $\FF_q$:

\begin{enumerate}[label=(\arabic*),start=5]
\item $\zeta(z) = \sum_{k \ge 0} z^k/[k]_q!$, from the Waldhausen $S_\bullet$-construction
on $\vect_q$; objectively, $q$-species with a version of the external product of
Joyal--Street \cite{JoyalStreet}.
\item $\zeta(z) = \sum_{k \ge 0} z^k/\norm{\Aut(\FF_q^k)}$, from the monoidal nerve of
$(\vect_q^{\iso},\oplus,0)$; objectively, the `Cauchy' product of $q$-species in the sense
of Morrison \cite{Morrison}.
\end{enumerate}

Their section coefficients are, respectively,
\[
  \qbin{n}{k}
  \qquad\text{and}\qquad
  \frac{\norm{\Aut(\FF_q^n)}}{\norm{\Aut(\FF_q^k)}\norm{\Aut(\FF_q^{n-k})}}
  \;=\; q^{k(n-k)}\qbin{n}{k}.
\]
The first counts subobjects with prescribed quotient (short exact sequences); the second
counts ordered direct-sum decompositions. The difference is $q^{k(n-k)} = \norm{U_P}$, the
order of the unipotent radical of the parabolic: the Waldhausen version induces from the
parabolic, the monoidal-nerve version from the Levi. GKT make the same point
structurally, observing that the map $\mathrm{Dec}_\bot \cat{M} \to \cat{W}$ is
\emph{not} an equivalence because ``taking complements'' has no inverse in the linear
case.

\begin{warning}[Prior art for \texttt{q-container.tex}]\label{warn:prior}
Proposition~\emph{(Day convolution over $\FF_q$)} and Proposition~\emph{(the Hall product
over $\FF_q$)} of the companion note \texttt{q-container.tex}, together with the
correction identifying the discrepancy as Levi-versus-parabolic, are the content of
\cite[\S{}Linear examples and the Waldhausen $S_\bullet$-construction]{GKT-comb}. The
$q$-binomial coalgebra there is credited further back to D\"ur \cite{Dur}. The companion
note's derivation is independent and its emphasis (containers, the Jackson derivative) is
not GKT's, but the two propositions should be cited to GKT and D\"ur rather than claimed.
\end{warning}

\begin{remark}[The two additive zetas differ by a Borel]
The ratio of the two normalisations is
\[
  \frac{\norm{\GL_n(\FF_q)}}{[n]_q!} \;=\; q^{\binom{n}{2}}(q-1)^n
  \;=\; \norm{U_B}\cdot\norm{T},
\]
verified symbolically for $n \le 8$. So the parabolic and Levi zeta functions differ by
the order of the Borel, factored into its unipotent and toral halves. Both halves recur
below, in different roles: $\norm{U_B}$ is the obstruction on the additive leg, $\norm{T}$
on the multiplicative one. That is the organising observation of this note.
\end{remark}

\subsection{The multiplicative leg, at \texorpdfstring{$q=1$}{q=1}}

For the classical multiplicative leg, GKT compute the monoidal nerve $\cat{A}$ of
$(\FinSet^{\iso}_{\neq\emptyset},\times,1)$, with section coefficients
\[
  \delta^m * \delta^n \;=\; \frac{(mn)!}{m!\,n!}\,\delta^{mn},
\]
and incidence algebra isomorphic to the exponential Dirichlet series via $\delta^m \mapsto
m^{-s}/m!$, with $\zeta = \sum_{k>0} k^{-s}/k!$. Baez \cite{Baez2025} pushes this further:
a scheme $X$ of finite type over $\ZZ$ determines a \emph{zeta species} $Z_X$, whose
Dirichlet series is the arithmetic zeta function of $X$; a $Z_X$-structure on a finite set
is a way of making it a semisimple commutative ring $k$ together with a $k$-point of $X$.
Finite fields are already present in that story --- but as the \emph{structures} being
counted, not as the base over which the categories are taken. It is a categorification of
$\prod_p \zeta_{X_p}$, not a deformation of $\zeta$.

\subsection{The gap}

Arranging the four combinations:

\begin{center}
\begin{tabular}{@{}lll@{}}
\toprule
 & additive ($\sqcup$, $\oplus$) & multiplicative ($\times$, $\otimes$) \\
\midrule
$\FinSet$ & species product, EGF & arithmetic product, Dirichlet \\
$\vect_q$ & Morrison / Joyal--Street, $q$-EGF & \textbf{---} \\
\bottomrule
\end{tabular}
\end{center}

The empty cell is the subject of the rest of this note. Searches of the arXiv corpus for
the conjunction of $q$-species, Dirichlet series, the Kronecker embedding $\GL_a \times
\GL_b \hookrightarrow \GL_{ab}$, and groupoid-cardinality zeta functions return nothing;
Morrison's $q$-species paper \cite{Morrison} treats subspaces, flags, direct-sum
decompositions and linear maps, all additive. We do not claim the cell is empty in the
literature at large, only that we did not find it.

\section{The multiplicative leg over \texorpdfstring{$\FF_q$}{Fq}}\label{sec:tensor}

\begin{definition}[$q$-arithmetic species]
Let $\vect_q^{\iso,\times} = \bigsqcup_{n\ge1} B\GL_n(\FF_q)$ be the groupoid of nonzero
finite-dimensional $\FF_q$-vector spaces and isomorphisms, monoidal under $\otimes$ with
unit $\FF_q$. A \emph{$q$-arithmetic species} is an object of
$\Grpd_{/\vect_q^{\iso,\times}}$, equivalently a functor $\vect_q^{\iso} \to \Grpd$. Their
convolution is Day convolution along $\otimes$:
\[
  (F \star G)(V) \;=\; \sum_{A \otimes B \,\cong\, V} F(A) \times G(B),
\]
the sum being over the homotopy fibre of $\otimes$ over $V$.
\end{definition}

Since $(\vect_q^{\iso},\otimes,\FF_q)$ is a monoidal groupoid, its monoidal nerve is a
Segal, hence decomposition, space, and GKT's machinery applies verbatim.

\begin{proposition}[Section coefficients]\label{prop:sect}
Writing $\delta^n$ for the name of an $n$-dimensional space,
\[
  \delta^a * \delta^b \;=\; c_{a,b}\,\delta^{ab},
  \qquad
  c_{a,b} \;=\; \frac{\norm{\GL_{ab}(\FF_q)}}{\norm{\GL_a(\FF_q)}\,\norm{\GL_b(\FF_q)}}.
\]
\end{proposition}

\begin{proof}
The homotopy fibre of $d_1$ over $V$ has objects $(A,B,\varphi)$ with $\varphi : A \otimes
B \xrightarrow{\ \sim\ } V$. Projecting to $\vect^{\iso}\times\vect^{\iso}$, the fibre over
$(A,B)$ is a set of cardinality $\norm{\Aut V}$, acted on by $\Aut A \times \Aut B$;
groupoid cardinality of the homotopy quotient is $\norm{\Aut V}/(\norm{\Aut A}\norm{\Aut
B})$ irrespective of freeness of the action. This is GKT's Corollary on section
coefficients for monoidal groupoids.
\end{proof}

\begin{remark}[The action is not free, and this is not a defect]
Unlike $S_a \times S_b \hookrightarrow S_{ab}$, the Kronecker embedding $\GL_a \times
\GL_b \to \GL_{ab}$ has kernel $\FF_q^\times$, embedded antidiagonally as $(\lambda,
\lambda^{-1})$. Hence the homotopy fibre is a groupoid with generic stabiliser
$\FF_q^\times$, and $c_{a,b}$ is in general \emph{not} an integer: the number of
$\GL_a\times\GL_b$-orbits is $(q-1)\,c_{a,b}$. Symbolic computation confirms that
$(q-1)c_{a,b} \in \ZZ[q]$ for all pairs tested. The multiplicative leg is intrinsically
stacky in a way the additive leg is not.
\end{remark}

\begin{proposition}[$q$-Dirichlet series]\label{prop:qdir}
The incidence algebra of the monoidal nerve of $(\vect_q^{\iso},\otimes,\FF_q)$ is
isomorphic to the algebra of formal Dirichlet series under the usual product, via
\[
  \Inc \;\longrightarrow\; \Bigl\{\textstyle\sum_{n>0} a_n n^{-s}\Bigr\},
  \qquad
  \delta^n \;\longmapsto\; \frac{n^{-s}}{\norm{\GL_n(\FF_q)}},
\]
under which the canonical zeta element $\zeta = \sum_{n>0}\delta^n$ maps to
\[
  \boxed{\;\zeta_q^{\otimes}(s) \;=\; \sum_{n \ge 1} \frac{n^{-s}}{\norm{\GL_n(\FF_q)}}\;}
\]
\end{proposition}

\begin{proof}
Setting $c_n = 1/\norm{\GL_n}$ we have $c_a c_b/c_{ab} = c_{a,b}$, which is exactly the
condition for $\delta^n \mapsto c_n n^{-s}$ to be multiplicative.
\end{proof}

This is the literal $q$-analogue of the Baez--Dolan zeta $\sum_{k>0} k^{-s}/k!$, with
$\norm{S_n}$ replaced by $\norm{\GL_n(\FF_q)}$ throughout. It is well defined, and it is
the honest output of the machinery. It is also, in two specific respects, wrong.

\section{The torus obstruction}\label{sec:torus}

\begin{lemma}[Why the additive leg works]\label{lem:cancel}
Write $\norm{\GL_n(\FF_q)} = \norm{U_B}\cdot\norm{T}\cdot[n]_q!$ with $\norm{U_B} =
q^{\binom n 2}$ and $\norm{T} = (q-1)^n$. In a convolution indexed by dimension, the toral
factor contributes
\[
  \frac{(q-1)^{\,\dim(A\otimes B)}}{(q-1)^{\dim A}(q-1)^{\dim B}}
  \;=\;
  (q-1)^{\,\dim(A\otimes B) - \dim A - \dim B}.
\]
For $\otimes = \oplus$ the exponent is $ (k + (n-k)) - k - (n-k) = 0$ identically, and the
torus cancels. For $\otimes$ the exponent is $ab - a - b$, which vanishes only for $(a,b)
\in \{(2,2)\}$ and for the degenerate cases involving $a$ or $b$ equal to $1$ with the
opposite sign.
\end{lemma}

\begin{proof}
Immediate; $\dim$ is additive for $\oplus$ and multiplicative for $\otimes$, while the
torus exponent is additive in both cases. Verified symbolically for all $a,b \le 5$.
\end{proof}

This single arithmetic fact accounts for all three pathologies of
Proposition~\ref{prop:qdir}.

\begin{corollary}[The unit is not $\delta^1$]
$c_{1,n} = \norm{\GL_n}/(\norm{\GL_1}\norm{\GL_n}) = 1/(q-1)$, independent of $n$. Hence
the convolution unit is $(q-1)\delta^1$, not $\delta^1$. Classically the unit object of
$(\FinSet,\times)$ is the singleton, with trivial automorphism group; over $\FF_q$ the
unit object is the line $\FF_q$, whose automorphism group is $\FF_q^\times$ of order
$q-1$. The multiplicative unit has a torus and cannot be rigidified away.
\end{corollary}

\begin{theorem}[No $\FF_1$ limit]\label{thm:nolimit}
$\lim_{q\to1} c_{a,b}$ exists and equals $(ab)!/(a!\,b!)$ if and only if $ab = a+b$. For
$ab > a+b$ the limit is $0$; for $ab < a+b$ it is $\infty$.
\end{theorem}

\begin{proof}
By Lemma~\ref{lem:cancel}, $c_{a,b} = (q-1)^{ab-a-b}\,q^{e}\,[ab]_q!/([a]_q![b]_q!)$ with
$e = \binom{ab}2 - \binom a2 - \binom b2$; the second and third factors have finite nonzero
limits $1$ and $(ab)!/(a!b!)$ respectively, so the limit is governed entirely by
$(q-1)^{ab-a-b}$.
\end{proof}

Explicitly, at $(a,b) = (2,3)$: $\lim_{q \to 1} c_{2,3} = 0$, against the classical
$6!/(2!3!) = 60$. The $\FF_1$ dictionary, which sends $\GL_n(\FF_q) \rightsquigarrow S_n$,
fails on the multiplicative leg. It does not fail because the dictionary is wrong; it
fails because the dictionary is $\norm{\GL_n(\FF_q)}/\norm{T} \to n!$, i.e.\ $\FF_1$-points
are monomial matrices modulo the torus, and only the additive leg lets the torus cancel by
itself.

\subsection{The rigidified normalisation, and what it costs}

\begin{definition}
Set $c_n^{\mathrm{rig}} = \norm{T}/\norm{\GL_n} = 1/(q^{\binom n2}[n]_q!)$, and
\[
  D_{a,b}(q) \;=\; \frac{c^{\mathrm{rig}}_a c^{\mathrm{rig}}_b}{c^{\mathrm{rig}}_{ab}}
  \;=\;
  q^{\binom{ab}{2}-\binom{a}{2}-\binom{b}{2}}\;\frac{[ab]_q!}{[a]_q!\,[b]_q!}.
\]
\end{definition}

\begin{theorem}\label{thm:D}
$D_{a,b}(q) \in \NN[q]$; it is a polynomial with nonnegative integer coefficients, of
degree $2\bigl(\binom{ab}2 - \binom a2 - \binom b2\bigr)$, symmetric in $a$ and $b$, with
$D_{1,n} = 1$ and $D_{a,b}(1) = (ab)!/(a!\,b!)$.
\end{theorem}

\begin{proof}[Verification]
Computed symbolically for all pairs in the table below; the degree formula and
nonnegativity were checked coefficientwise. We do not have a proof of nonnegativity in
general --- see Question~\ref{q:interp}.
\end{proof}

\begin{center}
\small
\begin{tabular}{@{}llrrr@{}}
\toprule
$(a,b)$ & $\deg D_{a,b}$ & $D_{a,b}(2)$ & $D_{a,b}(1)$ & $(ab)!/(a!b!)$ \\
\midrule
$(1,n)$ & $0$  & $1$ & $1$ & $1$ \\
$(2,2)$ & $8$  & $560$ & $6$ & $6$ \\
$(2,3)$ & $22$ & $19\,998\,720$ & $60$ & $60$ \\
$(2,4)$ & $42$ & $44\,213\,490\,155\,520$ & $840$ & $840$ \\
$(3,3)$ & $60$ & $24\,787\,851\,120\,790\,732\,800$ & $10\,080$ & $10\,080$ \\
$(2,5)$ & $68$ & $6\,107\,726\,516\,162\,836\,561\,920$ & $15\,120$ & $15\,120$ \\
$(3,4)$ & $114$ & $1.902\times10^{36}$ & $3\,326\,400$ & $3\,326\,400$ \\
\bottomrule
\end{tabular}
\end{center}

\begin{warning}[The repair is not objective]
$c^{\mathrm{rig}}_n$ is not the reciprocal of the cardinality of any automorphism group.
One might hope to realise it by projectivising --- replacing $B\GL_n$ by $B\PGL_n$, which
is natural given that $\FF_q^\times$ is exactly the kernel of the Kronecker embedding ---
but $\norm{\PGL_n} = \norm{\GL_n}/(q-1)$ removes one factor of $(q-1)$, not $n$ of them,
and the limit still fails for $ab - a - b \neq 1$. Dividing by the full diagonal torus is
not a quotient by a group acting on the groupoid. So
Theorem~\ref{thm:D} produces the right polynomial and the right $\FF_1$ limit at the cost
of the objective interpretation that motivated the construction.
\end{warning}

\begin{question}[Mai nafka minah]
Is there a monoidal groupoid, or more generally a decomposition space, whose section
coefficients are $D_{a,b}(q)$? Equivalently: is there a category of ``$\FF_q$-vector
spaces with a rigidified torus'' closed under $\otimes$? A positive answer would give the
multiplicative leg an objective footing; a proof of impossibility would establish that
Dirichlet convolution over $\FF_q$ is \emph{intrinsically} $\FF_1$-incompatible, which is
the more interesting outcome.
\end{question}

\begin{question}\label{q:interp}
$D_{a,b}(q)$ has nonnegative coefficients in every case tested. What does it count?
Classically $(ab)!/(a!b!)$ counts the ways of endowing a set of size $ab$ with an
$a \times b$ rectangle structure, divided by nothing --- it is
$\norm{S_{ab}}/(\norm{S_a}\norm{S_b})$, the index of the product action subgroup. The
$q$-analogue should count tensor decompositions of $\FF_q^{ab}$ modulo the torus, weighted
by some statistic; the $q$-exponent $\binom{ab}2 - \binom a2 - \binom b2$ suggests an
inversion-type statistic on pairs of flags.
\end{question}

\section{The \texorpdfstring{$q$}{q}-M\"obius function}

Being invertible ($c_1 \neq 0$), $\zeta_q^{\otimes}$ has a Dirichlet inverse
$\mu_q$. In the rigidified normalisation of Theorem~\ref{thm:D} the first values are:

\begin{center}
\small
\begin{tabular}{@{}lll@{}}
\toprule
$n$ & $\mu_q(n)$ & $\mu_q(n)\big|_{q\to1}$ \\
\midrule
$1$ & $1$ & $1$ \\
$2$ & $-\dfrac{1}{q(q+1)}$ & $-1/2$ \\[0.5em]
$3$ & $-\dfrac{1}{q^3(q+1)(q^2+q+1)}$ & $-1/6$ \\[0.5em]
$4$ & $\dfrac{q^8+q^7+2q^6+q^5+q^4-1}{q^6(q+1)^2(q^2+1)(q^2+q+1)}$ & $5/24$ \\[0.5em]
$5$ & $-\dfrac{1}{q^{10}(q+1)^2(q^2+1)(q^2+q+1)(q^4+q^3+q^2+q+1)}$ & $-1/120$ \\
\bottomrule
\end{tabular}
\end{center}

\begin{proposition}
For $p$ prime, $\mu_q(p) = -1/(q^{\binom p2}[p]_q!) = -c^{\mathrm{rig}}_p$, specialising to
$-1/p!$ at $q=1$.
\end{proposition}

\noindent
This is the $q$-analogue of the M\"obius function of the \emph{exponential} Dirichlet
zeta, not of Riemann's: at $q=1$ the values are $1, -\frac12, -\frac16, \frac5{24},
-\frac1{120}, \frac{119}{720}, \dots$, agreeing with $\bigl(\sum_k k^{-s}/k!\bigr)^{-1}$.
Note in particular that $\mu_q$ is not supported on squarefree integers and takes values
in $\QQ(q)$ rather than $\{0,\pm1\}$. That is already true classically for the fat
version; it is not a $q$-phenomenon.

\section{The other Dirichlet: \texorpdfstring{$\FF_q[t]$}{Fq[t]} and effective divisors}
\label{sec:functionfield}

There is a second, entirely different $\FF_q$-analogue of the second row of the table in
\S1, and it is the one arithmetic actually uses.

\begin{definition}
Let $\mathcal{M}_q$ be the multiplicative monoid of monic polynomials in $\FF_q[t]$,
viewed as a discrete monoidal category, with $\norm{f} = q^{\deg f}$. Day convolution on
$[\mathcal{M}_q, \Set]$ is
\[
  (F \star G)(f) \;=\; \coprod_{gh = f} F(g) \times G(h),
\]
decategorifying to the product of Dirichlet series $\sum_f a_f \norm{f}^{-s}$.
\end{definition}

Three things happen at once here, all of them absent from \S\ref{sec:tensor}.

\paragraph{The Euler product is a statement about Day convolution.}
$\mathcal{M}_q$ is the free commutative monoid on the monic irreducibles $P$. Day
convolution over a coproduct of monoids is the completed tensor product of the factors, so
\[
  [\mathcal{M}_q,\Set]_\star \;\simeq\; \widehat{\bigotimes_{P}}\; [\NN,\Set]_\star ,
\]
and the Euler product $\prod_P (1-\norm{P}^{-s})^{-1}$ is the decategorification of that
equivalence. It is not a theorem about zeta functions; it is the statement that Day
convolution respects free generation. Nothing of the sort holds in
\S\ref{sec:tensor}: $1/\norm{\GL_{mn}} \neq 1/(\norm{\GL_m}\norm{\GL_n})$ even for coprime
$m,n$, so $\zeta_q^{\otimes}$ has no Euler product.

\paragraph{Dirichlet collapses into Cauchy.}
$\deg : (\mathcal{M}_q,\times) \to (\NN,+)$ is strong monoidal, so pushing forward along it
carries Dirichlet convolution to Cauchy convolution. Writing $u = q^{-s}$, every Dirichlet
series over $\FF_q[t]$ is an ordinary power series in $u$, and
\[
  \zeta_{\FF_q[t]}(s) \;=\; \sum_{f} \norm{f}^{-s} \;=\; \sum_{n\ge0} q^n u^n
  \;=\; \frac{1}{1-qu}.
\]
Rationality of Weil zeta functions is, at this level, the assertion that the additive and
multiplicative legs are not independent over $\FF_q$ --- the grading fuses them. Over
$\ZZ$ there is no such grading, and $\zeta$ is not rational.

\paragraph{The primes are counted by the cyclotomic identity.}
Equating the Euler product with the closed form gives
\[
  \prod_{n\ge1}(1-u^n)^{-I(n)} \;=\; \frac{1}{1-qu},
  \qquad
  I(n) = \frac1n \sum_{d \mid n} \mu(d)\, q^{n/d},
\]
verified as formal power series to order $u^8$. $I(n)$ is simultaneously the number of
monic irreducibles of degree $n$, the number of primitive necklaces, and the dimension of
the degree-$n$ part of the free Lie algebra on $q$ generators. At $q=2$:
$2,1,2,3,6,9,\dots$

\begin{remark}
Everything in this section works verbatim for a smooth projective curve $X/\FF_q$ with
$\mathcal{M}_q$ replaced by the monoid of effective divisors, free commutative on the
closed points, giving $Z_X(u)$. Categorifying the \emph{fibres} rather than the index
monoid --- replacing divisors by coherent sheaves --- gives Kapranov's Hall algebra of
$X$.
\end{remark}

\section{Mai nafka minah}

Four constructions, all defensible as ``Day convolution over a finite field'':

\begin{center}
\small
\setlength{\tabcolsep}{4pt}
\begin{tabular}{@{}lcccc@{}}
\toprule
& Waldhausen & monoidal nerve & $\otimes$ & $\FF_q[t]$ \\
& (parabolic) & of $\oplus$ (Levi) & (this note) & (divisors) \\
\midrule
section coefficient & $\qbin nk$ & $q^{k(n-k)}\qbin nk$ & $\dfrac{\norm{\GL_{ab}}}{\norm{\GL_a}\norm{\GL_b}}$ & $1$ \\[0.8em]
zeta & $\sum \frac{z^n}{[n]_q!}$ & $\sum \frac{z^n}{\norm{\GL_n}}$ & $\sum \frac{n^{-s}}{\norm{\GL_n}}$ & $\frac{1}{1-qu}$ \\[0.6em]
$q$ lives in & automorphisms & automorphisms & automorphisms & index monoid \\
Euler product & no & no & no & \OK \\
$q \to 1$ limit & \OK & \OK & \NO\ (Thm~\ref{thm:nolimit}) & \OK \\
obstruction & --- & $\norm{U_P}$ & $\norm{T}$ & --- \\
in the literature & \cite{GKT-comb,JoyalStreet,Dur} & \cite{GKT-comb,Morrison} & --- & folklore \\
\bottomrule
\end{tabular}
\end{center}

The practical difference is sharp. If you want $q$-binomials, Hall--Littlewood and the
representation theory of $\GL_n(\FF_q)$, you want the first two columns, and the question
of which is settled by whether your induction is parabolic or Levi. If you want an Euler
product, a functional equation and the Riemann hypothesis, you want the fourth column, and
you should accept that the $q$ there is not deforming any automorphism group. The third
column is the literal transport of the classical Dirichlet construction, and
\S\ref{sec:torus} says it is obstructed.

\begin{remark}[Preserving the disagreement]
It is tempting to declare one of these ``the'' $q$-Dirichlet convolution. We think the
honest reading is that $(\NN_{>0},\times)$ does two different things at once --- it is a
lattice under divisibility, whose incidence algebra is Dirichlet convolution (Rota), and
it is a free commutative monoid on the primes, whose Day convolution is Dirichlet
convolution (Euler) --- and that these come apart over $\FF_q$. The lattice reading
deforms to the subspace lattice, hence to columns 1--2. The free-monoid reading deforms to
$\FF_q[t]$, hence to column 4. Column 3 is what you get by deforming the ambient rig
category instead of either reading, and it is the one that breaks.
\end{remark}

\section{What would do both}

The object that has $q$-deformed automorphisms \emph{and} an Euler product is the Hall
algebra of $\mathrm{Coh}(X)$ for $X/\FF_q$: locally at each closed point it is the Hall
algebra of finite modules over a discrete valuation ring, i.e.\ Hall--Littlewood; globally
its index monoid is the divisor monoid of \S\ref{sec:functionfield}. Kapranov--Schiffmann--
Vasserot's spherical Hall algebra of $\mathrm{Spec}(\ZZ)$ \cite{KSV} is the number-field
analogue. In the decomposition-space language, this is the Waldhausen
$S_\bullet$-construction of a category that is graded by an arithmetic monoid rather than
by $\NN$ --- both legs at once, which is precisely what \S\ref{sec:torus} shows cannot be
arranged inside $\vect_q$ alone.

\begin{question}[For the container programme]
The companion note \texttt{q-container.tex} takes positions in $\vect_q$. Replacing the
fibre category by finite $\FF_q[t]$-modules gives positions a genuine Dirichlet weight
$\norm{M} = q^{\ell(M)}$ alongside the additive $\oplus$; the Hall algebra of that
category is Hall--Littlewood and its groupoid cardinality $\sum_M 1/\norm{\Aut M}$ is
Cohen--Lenstra. Is there a container-level statement whose decategorification is the
functional equation?
\end{question}

\section{Summary of claims}

\begin{center}
\small
\begin{tabular}{@{}p{0.58\textwidth}>{\raggedright\arraybackslash}p{0.33\textwidth}@{}}
\toprule
claim & status \\
\midrule
Both additive legs over $\FF_q$, and the Levi/parabolic distinction, are in
\cite{GKT-comb,Dur,Morrison,JoyalStreet} & prior art, see Warning~\ref{warn:prior} \\
$\norm{\GL_n}/[n]_q! = q^{\binom n2}(q-1)^n = \norm{U_B}\norm{T}$ & \OK\ ($n\le8$) \\
Day convolution along $\otimes$ gives $c_{a,b} = \norm{\GL_{ab}}/(\norm{\GL_a}\norm{\GL_b})$ &
proved (Prop.~\ref{prop:sect}) \\
Incidence algebra $\cong$ Dirichlet series via $\delta^n \mapsto n^{-s}/\norm{\GL_n}$ &
proved (Prop.~\ref{prop:qdir}) \\
Convolution unit is $(q-1)\delta^1$ & proved \\
$c_{a,b}$ has no $q\to1$ limit unless $ab = a+b$ & proved (Thm~\ref{thm:nolimit}) \\
$D_{a,b}(q) \in \NN[q]$ with $D_{a,b}(1) = (ab)!/(a!b!)$ & \OK\ (tabulated pairs);
nonnegativity unproved in general \\
$\mu_q(p) = -1/(q^{\binom p2}[p]_q!)$ & \OK\ ($p \le 7$) \\
Cyclotomic identity / Euler product for $\FF_q[t]$ & \OK\ (to $u^8$); classical \\
No prior treatment of the $\otimes$ leg found & search result, not a proof \\
\bottomrule
\end{tabular}
\end{center}

All symbolic verification is in \texttt{q-dirichlet-verify.py}.

\begin{thebibliography}{99}

\bibitem{Baez2025} John C. Baez, \emph{Dirichlet species and arithmetic zeta functions},
arXiv:2502.01833.

\bibitem{BaezDolan} John C. Baez and James Dolan, \emph{Zeta functions},
\url{http://ncatlab.org/johnbaez/Zeta+functions} (2010).

\bibitem{CooperYoung} Benjamin Cooper and Matthew B. Young, \emph{Hall algebras via
2-Segal spaces}, arXiv:2409.19384.

\bibitem{Day} Brian Day, \emph{On closed categories of functors}, Reports of the Midwest
Category Seminar IV, Lecture Notes in Math. 137 (1970), 1--38.

\bibitem{Dur} Arne D\"ur, \emph{M\"obius functions, incidence algebras and power series
representations}, Lecture Notes in Math. 1202, Springer, 1986.

\bibitem{DyckerhoffKapranov} Tobias Dyckerhoff and Mikhail Kapranov, \emph{Higher Segal
spaces}, Lecture Notes in Math. 2244, Springer, 2019.

\bibitem{GKT-comb} Imma G\'alvez-Carrillo, Joachim Kock and Andrew Tonks,
\emph{Decomposition spaces in combinatorics}, arXiv:1612.09225.

\bibitem{GKT1} Imma G\'alvez-Carrillo, Joachim Kock and Andrew Tonks, \emph{Decomposition
spaces, incidence algebras and M\"obius inversion I: basic theory}, arXiv:1512.07573.

\bibitem{JoyalStreet} Andr\'e Joyal and Ross Street, \emph{The category of representations
of the general linear groups over a finite field}, J.\ Algebra 176 (1995), 908--946.

\bibitem{KSV} Mikhail Kapranov, Olivier Schiffmann and Eric Vasserot, \emph{The spherical
Hall algebra of $\mathrm{Spec}(\mathbb{Z})$}, arXiv:1202.4073.

\bibitem{MaiaMendez} Manuel Maia and Miguel M\'endez, \emph{On the arithmetic product of
combinatorial species}, arXiv:math/0503436.

\bibitem{Morrison} Kent E. Morrison, \emph{An introduction to $q$-species}, Electron.\ J.\
Combin.\ 12 (2005), arXiv:math/0512052.

\end{thebibliography}

\end{document}
